% !TeX program = xelatex
\documentclass[12pt, a4paper, titlepage, xelatex, ja=standard, jafont=hiragino-pron, twoside]{bxjsarticle}
\usepackage{xeCJK}
\xeCJKDeclareCharClass{CJK}{"203B}
\usepackage{amsmath, amssymb}
\usepackage{geometry}
\usepackage{longtable}
\usepackage{booktabs}
\usepackage{caption}
\usepackage{silence}
\WarningFilter{caption}{Unknown document class}
\WarningFilter{xeCJK}{Redefining CJKfamily}
\WarningFilter{relsize}{Failed to get list}

% -- Fonts --
\setmainfont{Times New Roman}

\geometry{reset, margin=2.5cm}

\title{ショートペーパー：仮想時間反復CCD法の性能評価と最適化戦略}
\author{}
\date{}

\begin{document}

\maketitle

\section*{1. 緒言 (Introduction)}
本研究の目的は、6次精度結合コンパクト差分（CCD）スキームを核とし、仮想時間反復（Pseudo-Time Iteration）を用いて支配方程式を解く手法の妥当性と効率性を検証することである。本手法の優位性を証明するため、従来の直接法や標準的な反復法に対し、精度・メモリ効率・計算時間の面から多角的な比較実験（Case 1-3）を実施する。

\section*{2. 実装の核：CCD演算子の一貫性}
すべての実験において、空間離散化は以下のCCD基本式をベースとする。未知数: 各格子点 $i$ における値 $\phi_i$、1階微分 $(\phi_x)_i$、2階微分 $(\phi_{xx})_i$。関係式: エルミート補完に基づくコンパクト形式。
$$\alpha (\phi_x)_{i-1} + (\phi_x)_i + \alpha (\phi_x)_{i+1} = \frac{a}{h}(\phi_{i+1} - \phi_{i-1}) + \dots$$
（※境界条件を含め、すべて $O(h^6)$ 精度で一貫させる）

\section*{3. 実験仕様 (Experimental Specifications)}

\subsection*{Case 1: ポアソン方程式（定常・精度検証）}
目的: 6次精度の完全な達成と、各ソルバーの収束特性の把握。

支配方程式: $\nabla^2 \phi = S(x, y)$

仮想時間導入: $\frac{\partial \phi}{\partial \tau} = \nabla^2 \phi - S$

比較指標:
\begin{itemize}
    \item 格子収束性 ($L_2, L_\infty$ 誤差の傾きが6であることの確認)
    \item ソルバー間の計算コスト比較（反復回数、メモリ、時間）
\end{itemize}

\subsection*{Case 2: 拡散方程式（非定常・安定性検証）}
目的: 物理時間進展における解法の安定性と効率性。

支配方程式: $\frac{\partial \phi}{\partial t} = \nu \nabla^2 \phi$

比較指標: 陽解法との比較による、物理時間刻み $\Delta t$ の制限（CFL条件）に対する堅牢性。

\subsection*{Case 3: 移流拡散方程式（非線形・界面模倣）}
目的: 移流卓越条件下での数値振動（Gibbs現象）の抑制能力。

支配方程式: $\frac{\partial \phi}{\partial t} + u \frac{\partial \phi}{\partial x} = \nu \nabla^2 \phi$

\section*{4. 比較対象ソルバー (Comparative Solvers)}
CCD離散化によって構成される連立方程式に対し、以下の手法を比較対象とする。

\begin{itemize}
    \item \textbf{直接法（LU分解）}: 基準解として使用。小規模格子での「正解」を得るが、大規模化に伴うメモリと時間の爆発を観測する。
    \item \textbf{陽解法（物理時間進展）}: 拡散数制限の厳しさを確認し、陰解法的アプローチの必要性を再認識する。
    \item \textbf{SOR法 (Successive Over-Relaxation)}: 古典的な反復法の代表。緩和係数 $\omega$ の最適化と、収束速度の限界を確認する。
    \item \textbf{BiCGSTAB法}: Krylov部分空間法の代表。CCD行列の非対称性に対する安定性と、前処理の有無による影響を確認する。
    \item \textbf{仮想時間反復法（提案手法）}: CCDのブロック三重対角構造を直接活かした反復。上記手法に対する優位性を検証する。
\end{itemize}

\section*{5. 評価軸 (Evaluation Metrics)}
純粋なアルゴリズム性能を測るため、以下の3点を主要な評価軸とする。

\begin{itemize}
    \item \textbf{反復回数 (Convergence Rate)}: 指定した収束残差（例：$10^{-10}$）に達するまでに必要な反復数。
    \item \textbf{メモリ使用量 (Memory Footprint)}: 巨大な係数行列を保持する手法（直接法、BiCGSTAB等）と、行列非保持あるいはブロック単位で処理する手法（仮想時間反復）の差。
    \item \textbf{総計算時間 (Total Computation Time)}: 同一精度に到達するまでの実実行時間。
\end{itemize}

\section*{6. 結論への道筋}
本実験を通じて、仮想時間反復CCD法が「高精度」と「計算資源の節約」を両立する次世代の数値流体解析の基盤（コア）になり得ることを、定量的かつ客観的なデータに基づいて証明する。

\end{document}